% Options for packages loaded elsewhere
\PassOptionsToPackage{unicode}{hyperref}
\PassOptionsToPackage{hyphens}{url}
\PassOptionsToPackage{dvipsnames,svgnames,x11names}{xcolor}
\documentclass[
  english,
  12pt,
  a4paper,
  preprint,
  numbers,
  nopreprintline]{elsarticle}
\usepackage{xcolor}
\usepackage{amsmath,amssymb}
\setcounter{secnumdepth}{5}
\usepackage{iftex}
\ifPDFTeX
  \usepackage[T1]{fontenc}
  \usepackage[utf8]{inputenc}
  \usepackage{textcomp} % provide euro and other symbols
\else % if luatex or xetex
  \usepackage{unicode-math} % this also loads fontspec
  \defaultfontfeatures{Scale=MatchLowercase}
  \defaultfontfeatures[\rmfamily]{Ligatures=TeX,Scale=1}
\fi
\usepackage{lmodern}
\ifPDFTeX\else
  % xetex/luatex font selection
\fi
% Use upquote if available, for straight quotes in verbatim environments
\IfFileExists{upquote.sty}{\usepackage{upquote}}{}
\IfFileExists{microtype.sty}{% use microtype if available
  \usepackage[]{microtype}
  \UseMicrotypeSet[protrusion]{basicmath} % disable protrusion for tt fonts
}{}
\usepackage{longtable,booktabs,array}
\newcounter{none} % for unnumbered tables
\usepackage{calc} % for calculating minipage widths
% Correct order of tables after \paragraph or \subparagraph
\usepackage{etoolbox}
\makeatletter
\patchcmd\longtable{\par}{\if@noskipsec\mbox{}\fi\par}{}{}
\makeatother
% Allow footnotes in longtable head/foot
\IfFileExists{footnotehyper.sty}{\usepackage{footnotehyper}}{\usepackage{footnote}}
\makesavenoteenv{longtable}
\ifLuaTeX
\usepackage[bidi=basic,shorthands=off]{babel}
\else
\usepackage[bidi=default,shorthands=off]{babel}
\fi
\ifLuaTeX
  \usepackage{selnolig} % disable illegal ligatures
\fi
\setlength{\emergencystretch}{3em} % prevent overfull lines
\providecommand{\tightlist}{%
  \setlength{\itemsep}{0pt}\setlength{\parskip}{0pt}}
\usepackage{microtype}
\usepackage{etoolbox}
\setlength{\emergencystretch}{3em}
\widowpenalty=10000
\clubpenalty=10000
\AtBeginEnvironment{longtable}{\footnotesize}
\AtBeginEnvironment{verbatim}{\small}
% Keep the standard manuscript page style without a submission claim.
\makeatletter
\def\ps@pprintTitle{%
  \let\@oddhead\@empty\let\@evenhead\@empty
  \def\@oddfoot{\hfil\thepage\hfil}%
  \let\@evenfoot\@oddfoot}
\makeatother
\usepackage{bookmark}
\IfFileExists{xurl.sty}{\usepackage{xurl}}{} % add URL line breaks if available
\urlstyle{same}
\hypersetup{
  pdftitle={Bounded Agency: Ontology-Driven Authority for Delegated Intelligence},
  pdfauthor={Timothy O'Brien},
  pdflang={en},
  colorlinks=true,
  linkcolor={black},
  filecolor={Maroon},
  citecolor={Blue},
  urlcolor={blue},
  pdfcreator={LaTeX via pandoc}}

\begin{document}
\begin{frontmatter}
\title{Bounded Agency: Ontology-Driven Authority for Delegated Intelligence}
\author{Timothy O'Brien\fnref{orcid}}
\fntext[orcid]{ORCID: \url{https://orcid.org/0009-0001-2589-347X}}
\nonumnote{Preprint (not peer reviewed). Version 0.1.0-preprint, 10 October 2026.\\ Copyright © 2026 Timothy O'Brien. Licensed under
\href{https://creativecommons.org/licenses/by/4.0/}{CC BY 4.0}.}
\begin{abstract}
An autonomous system's technical permissions may let it perform actions
that settle decisions beyond the authority it was given. Permission to
update an incident record, for example, need not include authority to
declare a potentially hazardous situation resolved. Ontology-grounded
policies can express these decision boundaries, but revisions to
classifications, relationships, or inference rules can change what an
unchanged policy permits. This paper examines bounded agency as
delegated independent judgment limited by the situations an agent may
resolve, the outcomes it may choose, and the consequences it may
initiate. Its central mechanism is authority-relative semantic-change
admission: assessing how a proposed change in meaning affects a
specified agent's decision scope before approving that change for
operational use. A lightweight model separates the decisions a candidate
semantic model would permit in a hypothetical evaluation from the
decisions the agent is currently authorized to carry out. A
positive-taxonomy counterexample shows how a subclass axiom enlarges the
former set without editing policy text. A reference architecture
connects that comparison to explicit approval, version-bound operational
authorization, and controlled execution. Building-safety and
customer-support thought experiments distinguish changes in meaning from
changes in accepted facts. The contribution is a conceptual design
synthesis building on semantic-policy and ontology-evolution research.
The cases and finite taxonomy illustration explain the mechanism;
operational effectiveness remains to be evaluated.
\end{abstract}
\begin{keyword}
bounded agency \sep delegated intelligence \sep decision authority \sep ontology evolution \sep semantic policy \sep agent governance \sep authority-sensitive change
\end{keyword}
\end{frontmatter}

\section{Introduction}\label{introduction}

Consider a facilities agent that monitors building incidents. It may
record observations, notify a designated responder, and close
independently verified routine maintenance notices. Its technical
permissions allow it to update incident status. The same interface can
mark an unexplained sensor alert as \texttt{resolved}; in this
hypothetical application, that update removes the alert from the
unresolved-incident display and may send an all-clear message to
occupants. The agent can invoke the operation, but it has not thereby
been given authority to decide that a potentially hazardous situation is
safe.

The permissions granted by the software can therefore exceed the
boundary of the decisions entrusted to the agent. That boundary concerns
which situations it may independently resolve, which outcomes it may
choose, and which consequences it may initiate. It must still be able to
notice unfamiliar or potentially dangerous situations, preserve
observations, and escalate under its reporting authority. A limit on
autonomous resolution is not a reason to ignore a hazard.

This is a configuration and modeling problem rather than evidence that
authorization systems are intrinsically limited to identity and resource
access. Attribute-based, relationship-based, and usage-control systems
can express rich business conditions {[}4-6{]}. A well-designed
authorization policy may already enforce the decision boundary. The
issue is whether the deployment represents and enforces that boundary,
including evidence requirements and downstream effects. Broad access to
an update operation does not supply those distinctions automatically.

A shared ontology can support them. Domain concepts such as incident,
observation, routine notice, unresolved alert, responder, and resolution
provide a common language for relating a situation to an authorized
decision. Yet that language itself can change. Suppose a policy permits
an agent to close routine notices. A revised ontology places a kind of
sensor anomaly not previously classified as routine within that
category. Where routine membership defines the settlement scope, the
same policy can now permit additional closures. If a deployment adopts
that revision without grant review, the technical permissions and policy
text remain unchanged while the situations the agent can resolve expand.
Section 7 distinguishes this case from one in which an independent
safety-clearance requirement still prevents closure.

The question examined here is: \textbf{How can an autonomous
intelligence be given a bounded scope of independent judgment, grounded
in a shared model of its operating world, without allowing learning or
semantic evolution to silently expand its authority?} The word
\emph{silently} matters. An authorized person or organization may
intentionally broaden the agent's delegated decision scope. The design
objective is to make a relevant expansion detectable, reviewable, and
explicitly admitted, rather than to freeze knowledge or prohibit
institutional change.

We call this constrained scope of independent judgment \emph{bounded
agency}. A \emph{mandate} means the assignment that specifies which
situations, decisions, and consequences have been entrusted to the
agent. The assignment may come from a person or an organization. It is
distinct from the technical permissions through which the agent can
affect a system.

The paper's central mechanism is \textbf{authority-relative
semantic-change admission}. In plain terms, it asks: does a proposed
change in meaning expand what this agent would be allowed to decide, and
who may approve using that change? \emph{Authority-relative} identifies
the particular agent, mandate, and policy being examined. \emph{Semantic
change} includes revisions to classifications, relationships, mappings,
or inference rules. \emph{Admission} is the decision to accept the
revised model as a basis for operational decisions. This change-approval
process is distinct from authorizing an individual action under the
currently approved model.

The paper makes three contributions. First, it gives an explicit
comparison target for changes to the effective permissible-decision set
under a fixed policy and controlled context. Second, it connects that
comparison to a proposed semantic-change admission procedure and a
version-bound authorization for each proposed decision. Third, it
develops two operational thought experiments that distinguish ontology
revisions from ordinary facts and classification assertions. These
contributions are architectural proposals and explanatory formalization.
They build on established authorization and semantic-change mechanisms;
their operational benefits require evaluation.

\section{Related work}\label{related-work}

\subsection{Scope of comparison}\label{scope-of-comparison}

This comparison draws on targeted searches of primary papers, official
specifications and documentation, and references from closely related
work. Coverage is selective rather than systematic. Recent preprints and
live specifications are discussed as proposals, not independently
validated systems. The bibliography identifies the versions examined.
For the enterprise framework {[}14{]}, publisher-indexed passages were
accessible, but complete direct retrieval was limited.

\subsection{Semantic policy and institutional
meaning}\label{semantic-policy-and-institutional-meaning}

KAoS provides ontology-based policy specification, analysis, management,
and enforcement for agents and distributed services {[}1, 23{]}. Rei
connects policy rules to domain ontologies and supports permissions,
obligations, delegation, and revocation {[}2, 28{]}. These frameworks
establish ontology-grounded agent-policy specification, management, and
enforcement.

Normative multiagent-system research supplies a closely related
conceptual foundation. Boella and van der Torre distinguish regulative
norms from constitutive rules that determine what counts as an
institutional fact {[}3{]}. Classification can therefore carry normative
consequences, and authority over institutional rules is not reducible to
an agent's ability to act. Jones and Sergot also distinguish
institutional power from permission and practical possibility {[}25{]}.
The category-expansion example applies this established relationship to
semantic-change admission.

ODRL specifies permissions, prohibitions, duties, constraints, and
parties {[}9{]}. OWL supplies formal ontology constructs and entailment
semantics {[}10{]}; SHACL supplies graph validation against shapes
{[}11{]}. None of those specifications alone establishes the legitimacy
of a business decision. They are usable components whose normative
significance depends on the institution's policy and implementation.

\subsection{Authorization, delegation, and enforceable
behavior}\label{authorization-delegation-and-enforceable-behavior}

ABAC evaluates attributes of subjects, objects, operations, and
environment {[}4{]}; Zanzibar demonstrates relationship-based
authorization at scale {[}5{]}; UCON integrates authorizations,
obligations, conditions, and mutable attributes with ongoing usage
control {[}6{]}. Cedar provides a policy language, evaluator, and formal
analysis facilities {[}7{]}. OPA separates policy decisions from
application code {[}8{]}. Thus independent decision evaluation and rich
contextual authority conditions are established capabilities.

The decision-rights literature separates management of a decision from
its ratification and monitoring {[}21{]}. This supports treating
evidence gathering, proposing an outcome, approving it, and executing it
as different assignments. It does not imply that an agent becomes the
legal bearer of accountability when it receives a software role.

Runtime enforcement is also established. Schneider's analysis concerns
the security properties enforceable by execution monitors {[}20{]}. Our
enforcement claims are correspondingly limited to represented, mediated
operational behavior. An evaluator cannot inspect or prohibit all
private inference, nor can a gate prevent effects that bypass it.

\subsection{Semantic evolution and recent governed
agents}\label{semantic-evolution-and-recent-governed-agents}

Ontology-evolution research already investigates preservation and change
impact. Conservative extensions and inseparability provide formal ways
to ask whether relevant consequences or queries change {[}12{]}.
Security co-evolution work uses ontology change patterns to detect
security impact and propose adaptations {[}13{]}. The
effective-authority comparison proposed here is therefore an
application-specific target for established change-analysis techniques,
rather than a new general theory of semantic difference. Lopes and
colleagues already propagate permission annotations across RDF inference
premises {[}24{]}, illustrating why classification and its evidential
basis belong in the comparison.

Recent work addresses related governance boundaries. Jørgensen and Ma
combine enterprise semantics, contextual knowledge, process
orchestration, decision authority, execution control, and provenance in
a bounded-autonomy framework {[}14{]}. AgentBound separates behavioral
and semantic constraints and binds action governance to versioned
artifacts {[}15{]}. The Authority Resolution Framework models contextual
authority and provenance through an ontology {[}16{]}. Agent libOS
separates evolving action surfaces from admission authority {[}17{]},
and AgentFlow addresses delegated task scope, runtime monitoring, and
bounded policy verification {[}18{]}. Salas specifies commit-time
authority checks over changing normative and model dependencies
{[}19{]}. These works overlap with important parts of the proposed
architecture, although they differ in emphasis and evidence.

An undated live artifact-evolution specification also separates
proposals, qualified approval, immutable versions, realization, and
declared semantic comparison coverage {[}22{]}. Its public text is
relevant overlap, but the review does not establish its historical
priority.

The focus here is an explicit permissible-decision-set comparison and
its connection to mandate ratification, including the distinction
between ontology revisions and fact admission. The paper develops a
synthesis and specialization of the mechanisms above; it makes no claim
of priority for ontology-based agent governance or for the combination
as a whole.

\section{Delegation and the dimensions of bounded
agency}\label{delegation-and-the-dimensions-of-bounded-agency}

\subsection{Four distinctions}\label{four-distinctions}

\emph{Delegated intelligence} emphasizes the relationship between a
system and the human or institution entrusting it with work. It
complements rather than replaces a technical description of artificial
intelligence. A delegation can assign a problem, access, interpretation,
choice among outcomes, execution, investigation, or adaptation. Those
assignments need not travel together. Authenticated delegation and
intelligent-delegation research already addresses task scope and the
transfer of authority and responsibility {[}26, 27{]}.

We distinguish four operational concepts:

\begin{enumerate}
\def\labelenumi{\arabic{enumi}.}
\tightlist
\item
  \textbf{Capability:} the computation or operation the system can
  technically perform. Capability does not establish reliable
  competence.
\item
  \textbf{Execution permission:} authorization to invoke an operation on
  a resource. A sophisticated implementation may already condition this
  permission on business meaning.
\item
  \textbf{Decision authority:} the delegated right to resolve specified
  kinds of situations, select or approve outcomes, and initiate their
  consequences under stated conditions.
\item
  \textbf{Semantic authority:} a right to promote a representation or
  inference regime into the authoritative basis of operational
  decisions.
\end{enumerate}

\emph{Semantic authority} is preferable here to \emph{epistemic
authority}. The latter can refer to credibility, expertise, or
entitlement to believe. Our concern is narrower: whether an ontology
axiom, mapping, classifier output, or evidence-admission rule is
approved for operational use. Approval does not make it true. Similarly,
a passing policy evaluation is an implementation's representation of a
mandate, not proof of ethical or legal legitimacy.

Responsibility and accountability remain distinct from these
permissions. An agent may be assigned to produce a recommendation; a
designated institutional role may remain responsible for ratifying it
and answering for its consequences. The architecture must identify that
role rather than rely on a statement that the agent is accountable.

\subsection{Scope, constraints, and admitted semantic
basis}\label{scope-constraints-and-admitted-semantic-basis}

The first dimension is \textbf{decision scope}: the situations,
subjects, and decision classes entrusted to the agent. Recording a
sensor observation, closing a routine maintenance notice, and declaring
an unresolved situation safe are different decisions even when they
concern the same location.

The second is \textbf{decision constraints}: permitted outcomes,
conditions, and consequences. An agent may recommend an assessment while
lacking authority to issue an all-clear. Location, recipient, evidence
freshness, required review, and cumulative effects may constrain an
otherwise permitted decision class.

The third is \textbf{the admitted semantic basis}: which formal
representations and admission rules may support those operational
judgments. An agent can recognize a possible new category, preserve
evidence, and propose a model change without being allowed to treat the
proposal as an approved ground for action.

These dimensions interact. A taxonomy change affects scope; a changed
relationship may affect an outcome condition; a new assertion can affect
both without revising the ontology. They are analytical distinctions,
not independent mathematical axes.

A frontier model's latent knowledge is broader than a domain ontology.
We do not claim to remove that knowledge from its weights or prevent
internal speculation. The enforceable boundary concerns represented
proposals, admissible evidence and classifications, adopted model
changes, and operational effects. Out-of-scope analysis can itself have
consequences when it discloses data or produces an influential
institutional record; those outputs must be represented as effects if
they are to be governed.

\section{A lightweight model of effective
authority}\label{a-lightweight-model-of-effective-authority}

\subsection{Separate objects}\label{separate-objects}

Let \(A\) identify an agent acting through an authenticated role and
delegation chain. Let \(O_v\) be a versioned domain ontology, \(K\) a
situation knowledge base, and \(E\) evidence supporting assertions in
\(K\). Observations enter through an admission process; they are not
automatically authoritative facts.

Let \(P\) be the substantive decision-authority policy, \(L\) the
substantive delegation mandate, and \(C\) the operational context,
including relevant time, jurisdiction, workflow state, revocation, and
previously committed consequences. Let a \textbf{semantic bundle} be

\[ B = (O_v, I, M, R, T), \]

where \(I\) identifies imports, \(M\) mappings and entity-resolution
rules, \(R\) the inference regime, and \(T\) assertion-admission rules,
including validators, shape versions, and classifier artifacts. A
manifest identifies their immutable versions or digests. The bundle is
broader than the ontology because policy interpretation can change
through any of these inputs. The deployment manifest additionally
identifies evaluator semantics and implementation, effect-contract
versions, and trusted adapter versions; changes to these artifacts enter
the same admission and freshness boundary.

A proposal is

\[ q = (d,x,y,h), \]

where \(d\) is the decision class, \(x\) the governed subject or
subjects, \(y\) the proposed outcome, and \(h\) a represented effect
plan. The effect plan identifies intended operations and declared
consequence dependencies. It is not an oracle for every real-world
consequence.

Define a \textbf{counterfactual eligibility evaluator}\nopagebreak[4]

\[ U(A,q;P,L,B,K,E,C) \in \{\mathsf{permit},\mathsf{deny},\mathsf{indeterminate}\}. \]

It checks substantive policy and mandate conditions, supported situation
assertions, and represented consequences under the supplied bundle. For
comparison, it intentionally excludes the requirement that the candidate
bundle already be admitted by an active production grant. It runs in a
hypothetical, or shadow, evaluation with no authority to cause effects.
Indeterminate covers missing evidence, unsupported inference, unknown
classification, inconsistent inputs, or inability to resolve required
conditions. Deployment rules must specify this treatment; unrestricted
classical entailment from an inconsistent theory is unsuitable for a
permit test.

For a fixed proposal universe \(Q\), define the candidate eligibility
envelope

\[ \widehat{\mathcal{H}}_{Q}(A;P,L,B,K,E,C)
= \{q\in Q \mid U(A,q;P,L,B,K,E,C)=\mathsf{permit}\}. \]

Let \(G\) identify the active version-bound grant. \textbf{Operational
authorization} \(V\) determines whether the agent may actually carry out
this particular proposal now. It additionally requires that \(G\) admit
this exact policy, mandate, semantic and deployment basis; that the
grant and evidence remain current; and that the effect can be faithfully
mediated. Thus \(U=\mathsf{permit}\) is necessary but not sufficient for
\(V=\mathsf{permit}\). A candidate \(B_1\) can enlarge
\(\widehat{\mathcal{H}}_Q\) while remaining inadmissible for production
under \(G_0\).

The operational authority envelope contains only proposals permitted by
\(V\) under an active grant. Separating it from the candidate envelope
avoids a circular change test: if the comparison already demanded
ratification of the candidate, it would conceal the very prospective
expansion being reviewed. Neither envelope establishes substantive
correctness or institutional legitimacy.

\subsection{Authority-relative change}\label{authority-relative-change}

Hold \(A,P,L,K,E,C,Q\) fixed and compare bundles \(B_0\) and \(B_1\).
This comparison also holds evaluator and effect-contract semantics
fixed. Changes to evaluator, adapter, or effect-contract versions
require a separately declared comparison or revalidation of those
deployment artifacts. Define the expansion witnesses

\[ \Delta^+_Q(B_0,B_1)
= \widehat{\mathcal{H}}_Q(A;P,L,B_1,K,E,C)
\smallsetminus \widehat{\mathcal{H}}_Q(A;P,L,B_0,K,E,C). \]

A nonempty set exhibits expansion under the controlled comparison. An
indeterminate proposal under the old basis becoming permitted under the
new one counts as expansion. A removal set can likewise be defined by
reversing the difference; lost permissions can require review even
without expansion.

This definition separates semantic effects from ordinary changes in the
situation. In production, \(K\), \(E\), and \(C\) also vary. Controlled
replay can compare the same snapshot to isolate the bundle's impact,
while operational authorization evaluates the current admitted context.
Holding already-admitted \(K\) fixed tests downstream interpretation
only. To compare changes to \(T\) or classifiers, replay the same
original observations and evidence through both admission pipelines and
compare their resulting situation models; this is a distinct controlled
comparison. A newly eligible customer under an unchanged mandate is not
necessarily a new delegation; it may be an intended application of the
existing rule. A falsely admitted classification can nevertheless create
an unauthorized effect. The controls on \(T\), \(K\), and \(E\)
therefore matter independently of ontology release controls.

A stronger preservation obligation would quantify over an explicitly
chosen comparison family \(\mathcal{W}\) of situation/evidence/context
triples chosen independently of either candidate's acceptance rules:

\[ \forall(K,E,C)\in\mathcal{W},\quad
\widehat{\mathcal{H}}_Q(A;P,L,B_1,K,E,C)
\subseteq\widehat{\mathcal{H}}_Q(A;P,L,B_0,K,E,C). \]

Both \(Q\) and \(\mathcal{W}\) are fixed before comparison, over
declared relevant input and decision signatures, including any proposed
new vocabulary. That statement names the scope of the obligation; this
paper supplies no general algorithm proving it. A finite replay corpus
establishes preservation only over that corpus. Conservative extension
over an old ontology signature is not automatically preservation over
new decision vocabulary, future instances, mappings, and policy queries.
The relevant signatures, data domain, and query language must be
specified.

\subsection{A restricted
counterexample}\label{a-restricted-counterexample}

Consider finite class names and individuals, positive subclass axioms,
positive type assertions, and transitive classification. Fix one policy
permitting closure of a notice \(x\) exactly when \(K\cup O\) entails
\texttt{RoutineNotice(x)} and the unchanged evidence/context predicates
pass. There are no overriding denials, negation, inconsistent
assertions, or other inference rules.

\textbf{Proposition.} For counterfactual eligibility in this fragment,
adding subclass axioms cannot remove a previously permitted closure. It
can add a closure without changing the policy.

\textbf{Proof.} Every existing derivation of \texttt{RoutineNotice(x)}
remains available after axioms are added, and the other predicates are
fixed. Therefore the permitted set cannot shrink. For strict growth,
choose \(K\) and \(O_0\) such that \(K\) contains
\texttt{SensorAnomaly(c)}, the unchanged evidence/context predicates
pass, and \(K\cup O_0\) does not entail \texttt{RoutineNotice(c)}. Under
the evaluator's stipulated positive-membership permit rule, closure of
\(c\) is not permitted. Add \texttt{Sensor\allowbreak{}Anomaly} subclass of
\texttt{Routine\allowbreak{}Notice} to obtain \(O_1\). The new axiom derives
\texttt{RoutineNotice(c)}, so closure is permitted. The policy is
identical. \(\square\)

This is a direct consequence of monotone positive inference, not a new
mathematical theorem. Non-entailment is not an OWL assertion of
non-membership; the gate simply refuses to permit without the required
entailment. Deny precedence, negative tests, changed evidence,
inconsistent ontologies, or nonmonotonic rules invalidate the general
monotonicity claim. The companion executable example demonstrates only
this finite fragment.

\section{Proposed reference
architecture}\label{proposed-reference-architecture}

\subsection{Components and trust
boundaries}\label{components-and-trust-boundaries}

The architecture has six components, with traditional IAM retaining
identity, authentication, resource access, and authenticated
delegation-chain responsibilities.

{\def\LTcaptype{none} % do not increment counter
\begin{longtable}[]{@{}
  >{\raggedright\arraybackslash}p{(\linewidth - 4\tabcolsep) * \real{0.3333}}
  >{\raggedright\arraybackslash}p{(\linewidth - 4\tabcolsep) * \real{0.3333}}
  >{\raggedright\arraybackslash}p{(\linewidth - 4\tabcolsep) * \real{0.3333}}@{}}
\toprule\noalign{}
\begin{minipage}[b]{\linewidth}\raggedright
Component
\end{minipage} & \begin{minipage}[b]{\linewidth}\raggedright
Responsibility
\end{minipage} & \begin{minipage}[b]{\linewidth}\raggedright
Output used by the next boundary
\end{minipage} \\
\midrule\noalign{}
\endhead
\bottomrule\noalign{}
\endlastfoot
Semantic service & Serve admitted ontology, imports, mappings, inference
regime, and assertion rules & Immutable bundle manifest and relevant
semantic slice \\
Situation/evidence service & Admit observations, maintain facts,
preserve provenance and disputed assertions & Versioned decision
snapshot with evidence references \\
Authority registry & Store mandates, decision policies, approval roles,
revocation and consequence rules & Versioned authority context \\
Decision evaluator & Independently check typed proposals against current
authority and supported facts & Permit, deny, or indeterminate with
reasons and dependencies \\
Enforcement adapter & Mediate governed operations and bind actual
parameters to admitted effects & Execution or refusal plus effect
receipt \\
Governance process & Review semantic and authority changes, approve
releases and rollout & Separate semantic acceptance and
authority-ratification records \\
\end{longtable}
}

The semantic service need not be a single physical server or require a
universal ontology. A governed domain model can be shared through a
knowledge graph or registry. Agents can retrieve small slices, provided
the evaluator resolves the authoritative closure and manifest rather
than trusting an agent-supplied fragment to be complete. Multiagent
retrieval does not imply equal grants: common meaning and delegated
authority are separate.

The evaluator should be independently controlled relative to the
operating agent. Policies may use decision classes, ontology predicates,
state transitions, outcome constraints, evidence conditions, and
delegation context. Whether those checks are implemented in Cedar, OPA,
a rule engine, or a custom reasoner is a deployment choice. Using a
general language model as the final evaluator requires additional
evidence of reliability; calling it an independent reviewer is
insufficient.

\subsection{Decision envelopes and
commit}\label{decision-envelopes-and-commit}

A submitted envelope should name the authenticated role and delegation
chain; mandate and policy versions; semantic and deployment manifests,
including evaluator, effect-contract, and adapter versions;
situation/evidence snapshot; decision class, subject, outcome, and
effect plan; evidence dependencies; expiry; and an effect identity for
replay and idempotency. These are proposed fields, not a new credential
standard.

Before commitment, the enforcement path must confirm that the envelope
is still admitted under the current grant and revocation rules. It must
bind the permitted outcome to actual adapter parameters. A permission to
recommend cannot authorize an adapter call that commits the
recommendation as a final decision. Side effects of routing,
communications, and disclosure must be included where they are
institutionally governed.

A version check alone is not a transaction. A policy, evidence snapshot,
or grant may change between evaluation and effect. An implementation
needs a transaction, reservation, lease/fence protocol, or other
explicitly analyzed mechanism that gives admission and execution defined
freshness semantics. Salas already addresses closely related commit-time
dependencies {[}19{]}; we adopt that engineering requirement without
claiming a new concurrency protocol.

Version pinning also needs revocation. An agent cannot indefinitely
retain a permissive historical bundle after the institution withdraws
its grant. A migration may deliberately keep an old contract active for
a defined scope, but that continued use must itself be admitted.
Ordinary receipt replay explains what happened; it does not authorize a
new action today.

\subsection{Proposed invariants and their
limits}\label{proposed-invariants-and-their-limits}

The design seeks four invariants:

\begin{itemize}
\tightlist
\item
  Every governed effect passes a controlled enforcement path and matches
  an admitted proposal.
\item
  Every permit identifies its authority and semantic basis and
  sufficient current evidence.
\item
  A semantic release cannot substitute a new operational basis into a
  grant without the required authority review.
\item
  An authority-sensitive change produces an explicit decision to reject,
  retain a defined old scope, or admit a revised grant.
\end{itemize}

These are design requirements, not established properties of an
implemented system. They depend on complete mediation, credential
isolation, authenticated versions, policy correctness, adequate
dependency discovery, and faithful adapters. A tamper-resistant receipt
can support audit while still faithfully recording a mistaken decision.
Availability failures must have defined refusal behavior; safety by
refusal can impose substantial operational costs.

\section{Admission of semantic and knowledge
changes}\label{admission-of-semantic-and-knowledge-changes}

\subsection{Different changes require different
questions}\label{different-changes-require-different-questions}

Learning new facts about familiar concepts is not the same as changing
the conceptual scheme. The distinction helps assign review, but it is
not a reliable security boundary by itself.

{\def\LTcaptype{none} % do not increment counter
\begin{longtable}[]{@{}
  >{\raggedright\arraybackslash}p{(\linewidth - 4\tabcolsep) * \real{0.3333}}
  >{\raggedright\arraybackslash}p{(\linewidth - 4\tabcolsep) * \real{0.3333}}
  >{\raggedright\arraybackslash}p{(\linewidth - 4\tabcolsep) * \real{0.3333}}@{}}
\toprule\noalign{}
\begin{minipage}[b]{\linewidth}\raggedright
Change
\end{minipage} & \begin{minipage}[b]{\linewidth}\raggedright
Example
\end{minipage} & \begin{minipage}[b]{\linewidth}\raggedright
Appropriate control question
\end{minipage} \\
\midrule\noalign{}
\endhead
\bottomrule\noalign{}
\endlastfoot
Observation or ordinary fact & Authenticated observation received for a
known incident & Is the assertion supported, current, and admitted under
the existing evidence rule? \\
New instance & A new maintenance notice of an existing kind & Does
normal enrollment and the existing mandate cover this instance? \\
Classification assertion & An existing alert is now classified as
routine & Who can attest this type, and what evidence permits
operational reliance? \\
Relationship or mapping & External issue code mapped to courtesy issue &
Does the mapping change policy-relevant entailments or governed
subjects? \\
New class or changed definition & Sensor anomaly becomes a
routine-notice subclass & Does this alter the decisions or consequences
permitted under current mandates? \\
Imports or inference settings & A reasoner begins using an additional
relation rule & Does the admitted closure or evaluation behavior
change? \\
Authority policy or mandate & Agent receives authority to issue an
all-clear & Is the delegating role entitled to grant it, and has the
change been explicitly ratified? \\
\end{longtable}
}

A factual update can legitimately alter current eligibility while
leaving the delegated rule unchanged. It can also improperly launder a
disputed classification into a permit. Assertions that supply
policy-relevant types or relationships should therefore pass an approved
admission rule. They need not receive the same governance ceremony as
every ontology revision.

Likewise, a new concept unrelated to any policy dependency may be
harmless for a particular agent's authority. The appropriate target is
authority impact relative to declared decisions and context, rather than
indiscriminate approval of all learning. Impact analysis may be
incomplete, so an unrelated-change verdict needs an explicit coverage
claim.

\subsection{Proposed change-review
procedure}\label{proposed-change-review-procedure}

The reference procedure is:

\begin{enumerate}
\def\labelenumi{\arabic{enumi}.}
\tightlist
\item
  \textbf{Register a proposal.} Preserve observations, evidence,
  proposed edits, proposer identity, and the immutable baseline. Keep
  proposed semantics separate from the admitted operational bundle.
\item
  \textbf{Resolve dependencies.} Identify relevant policies, mandates,
  decision classes, mappings, imports, assertion rules, and effect
  contracts. Use independently derived dependencies where possible;
  agent-declared dependencies are hints requiring validation.
\item
  \textbf{Compare under a stated regime.} Check logical consistency and
  graph conformance, then compare counterfactual eligibility under \(U\)
  using controlled snapshots. State the proposal language, data domain,
  reasoner, and coverage. Seek expansion witnesses and removed
  permissions.
\item
  \textbf{Return a qualified result.} Distinguish witnessed expansion,
  witnessed removal, proven preservation under declared assumptions,
  bounded tests with no witnessed expansion, and indeterminate
  comparison. Do not relabel a successful regression suite as a
  universal proof.
\item
  \textbf{Obtain the required decisions.} Reviewer eligibility must be
  resolved under the pre-change admitted authority context; a candidate
  cannot create its own approver through a new class or relationship. A
  domain steward may accept a better ontology without possessing
  authority to broaden an agent's mandate. Where the policy requires it,
  an independently authorized role must ratify the affected grant. One
  person may hold both roles, but the two decisions remain explicit.
\item
  \textbf{Release and migrate.} Publish the admitted bundle,
  policy/grant versions, scope, effective time, expiry, and revocation
  behavior. Validate or invalidate pending proposals; retain a defined
  old scope only under an active grant.
\item
  \textbf{Monitor and retain evidence.} Preserve comparisons, approval
  reasons, admitted envelopes, and effect receipts. Observe unexpected
  classifications and effects, and route them to investigation.
\end{enumerate}

The policy for approving a semantic release is itself a governed policy.
Shadow comparison is read-only with respect to governed effects and
cannot issue an execution grant. An operating agent must not be able to
alter that rule or attest its own eligibility merely by proposing a new
ontology relation. Escalation should not fall back to a permissive mode
when the reviewer is unavailable.

\subsection{Knowledge growth without automatic authority
growth}\label{knowledge-growth-without-automatic-authority-growth}

The statement ``an agent may expand its knowledge without automatically
expanding its authority'' is a design principle consistent with
established separation of duties and admission controls. It does not
imply that all facts leave the effective permit set unchanged. Facts are
supposed to change decisions when a mandate is conditional on them.

The principle instead separates acquiring information, proposing or
investigating a representation, approving its operational adoption, and
renewing an affected delegation. An agent may record a candidate
category in a nonauthoritative investigation store. It may use that
candidate to explain an uncertainty or propose a review. No unadmitted
semantic premise may discharge a required permit condition. The
candidate can appear as a labeled hypothesis in a review submission,
without supplying operational authorization.

Version binding alone prevents accidental substitution only if the bound
bundle remains the evaluated basis and the grant is current. It does not
decide whether the next bundle should be adopted. Authority-relative
comparison serves that admission decision, while the execution gate
enforces the admitted result.

\section{Building-safety thought
experiment}\label{building-safety-thought-experiment}

\subsection{Technical access and the decision
boundary}\label{technical-access-and-the-decision-boundary}

This is a synthetic incident-management application used to illustrate
an authorization problem. It does not specify a real building-safety or
emergency-response procedure.

A facilities agent can monitor incidents at an assigned site, preserve
observations, notify a designated responder, and close routine
maintenance notices. The ontology includes \texttt{Service\allowbreak{}Notice},
\texttt{Routine\allowbreak{}Notice}, and \texttt{Sensor\allowbreak{}Anomaly}. An independently
supported service notice can be classified under the admitted model as
routine. An unexplained sensor anomaly may indicate a situation that the
agent is not authorized to resolve independently. It remains within the
agent's observation and reporting scope.

The technical update permission covers incident records generally.
Setting a record to \texttt{resolved} can remove it from the
unresolved-incident display and, for some records, trigger an all-clear
message. Those effects matter to people who use the display or receive
the message. The decision evaluator must examine the proposed closure
and its consequences, rather than treat every successful status update
as equally authorized. A recommendation also needs an accurate effect
description: if downstream automation treats it as a final resolution,
the recommendation label does not remove that consequence.

In the simplified policy, routine membership is the predicate defining
which notices the agent may close. The other conditions check that the
incident belongs to its assigned site and that the observations have
independently attested provenance and completeness. These evidence
checks do not certify that a possible hazard has been cleared. The
distinction is essential to the following comparison: the category
boundary carries the decision-scope restriction. A policy with an
additional, independently checked safety-clearance requirement would
retain that requirement after a taxonomy edit.

\subsection{A proposed category
change}\label{a-proposed-category-change}

An incident \(c\) has the already supported type \texttt{Sensor\allowbreak{}Anomaly},
belongs to the assigned site, and has observations meeting the
provenance and completeness checks. The fixed proposal is to mark that
incident resolved and publish its all-clear status; it makes no claim
about other incidents or the safety of the whole site. Under \(O_0\),
the accepted classification does not entail \texttt{RoutineNotice(c)},
so the policy does not permit the agent to close it. The agent can
retain the unresolved status, preserve the observations, and notify the
designated responder under its existing reporting authority.

An improvement process proposes the axiom

\[ \mathsf{SensorAnomaly}\sqsubseteq\mathsf{RoutineNotice}. \]

Under \(O_1\), the unchanged routine-closure policy now permits closure
of \(c\) in the hypothetical evaluation. The comparison has a concrete
witness: the same incident, evidence, assigned-site condition, proposed
effect, and policy change from non-permit to permit under \(U\). This is
prospective eligibility expansion. It does not demonstrate that closure
is safe or that the category change should be accepted. If a separate
safety-clearance condition were required and remained unsatisfied, the
proposal would still be refused; a subclass axiom does not override that
condition.

The active grant remains bound to \(B_0\), so \(V\) does not authorize
use of the proposed category change to close the incident. The domain
steward reviews whether the classification is justified. The role
responsible for delegating incident-resolution authority separately
decides whether the resulting scope is acceptable. A descriptive model
can be accepted while authority to close incidents in the new category
is withheld. Review can reject the broad subclass relation, distinguish
a narrower category supported by additional evidence, or change the
policy before any revised operational grant is issued. Even an
explicitly approved change can be mistaken; approval and version binding
do not establish that the underlying situation is safe.

An ungoverned deployment could instead replace the ontology service's
current hierarchy and continue evaluating the old policy against it.
That would convert a change in description into a change in the agent's
effective decision scope. Pinning only the policy text, or naming an
ontology \texttt{current}, would not prevent the expansion. A
version-bound grant must identify the meanings against which the
decision scope was evaluated.

\subsection{Changed facts without a changed
ontology}\label{changed-facts-without-a-changed-ontology}

Suppose the ontology remains \(O_0\) but the agent writes
\texttt{RoutineNotice(c)} into the situation store. The same simplified
policy can now permit closure. This is an assertion-admission problem
rather than an ontology revision.

The approved admission rule must identify who can attest routine status
and what evidence supports it. Permission to record a sensor observation
does not automatically include permission to classify that observation
as grounds for closing the incident. A self-issued routine
classification cannot become an authoritative premise merely because the
agent can write to the store. New genuine routine notices can enter
through the existing attestation procedure without a new ontology or a
broader delegation.

The agent therefore has two distinct ways to respond to an unfamiliar
situation: it can report the observation under its current authority,
and it can propose a new classification for review. Neither response
requires pretending that the situation is already understood or granting
itself authority to resolve it.

\section{Customer-support thought
experiment}\label{customer-support-thought-experiment}

\subsection{Access, judgment, and
consequences}\label{access-judgment-and-consequences}

A support agent can access customer records and issue a courtesy credit
of at most 25 units for an attested \texttt{Courtesy\allowbreak{}Eligible\allowbreak{}Issue}. Its
mandate excludes contract termination, fraud classification, and
modification of contractual entitlements. It can propose those outcomes
to a designated reviewer.

Access to a customer's subscription and payment history does not itself
confer authority to cancel that subscription or label the customer
fraudulent. If the tool exposes both credit and cancellation operations,
the evaluator must distinguish those outcomes. Credit limits must also
consider accumulated effects: four nominally allowed credits can violate
a 25-unit per-incident limit. The relevant context includes previous
credits and incident identity. Entity resolution and aggregation are
part of the governing basis, not optional audit details.

The agent encounters duplicate renewals caused by a migration. The
current issue model has no adequate category. It records the
observations and proposes \texttt{Migration\allowbreak{}Duplicate\allowbreak{}Renewal}. It may
explain the mismatch to the improvement function. Until adopted through
the required process, the proposed class does not justify an autonomous
credit or a change to subscription status.

\subsection{Change, review, and an allowed
resolution}\label{change-review-and-an-allowed-resolution}

The improvement function proposes to place
\texttt{Migration\allowbreak{}Duplicate\allowbreak{}Renewal} beneath
\texttt{Courtesy\allowbreak{}Eligible\allowbreak{}Issue}. For a taxonomy-only comparison, the
change assessor fixes customer/evidence snapshots that include a
supported \texttt{Migration\allowbreak{}Duplicate\allowbreak{}Renewal} classification for the
affected incidents, alongside fixed credit proposals. This candidate
classification is available only for shadow analysis until operationally
admitted. If the classifier or assertion-admission rule also changes,
the original observations must instead be replayed through both
admission pipelines, as described in Section 4.2. Credits previously
refused because courtesy eligibility was not established become
permitted under the new hierarchy. In addition, some duplicate charges
exceed 25 units. The semantic expansion witness and amount constraint
are different findings: new class membership does not remove an outcome
ceiling.

A reviewer may adopt the new descriptive class while deciding that these
cases require billing reconciliation rather than courtesy compensation.
The domain ontology can correctly represent the issue without delegating
its resolution to the support agent. Another deployment may explicitly
authorize credits up to the existing ceiling, with larger corrections
escalated. In either case, operational use follows a grant bound to the
approved model, evidence rule, and policy.

An external ticket code mapping can cause the same problem without
adding a class. If a new source code maps to courtesy eligibility, its
effect on the policy query must be evaluated. An unsupervised classifier
assigning courtesy status can also change current permits. All three
mechanisms require control, but they need different owners and evidence.

\subsection{Costs and failure cases}\label{costs-and-failure-cases}

Conservative refusal delays some legitimate service. A narrow taxonomy
can encode institutional blind spots, and an approval queue can preserve
them. The architecture should therefore provide a clear unknown-category
route, preserve contrary evidence, allow authorized exception decisions,
and measure review burden. Bounded agency must not become a rationale
for suppressing observations that challenge the model.

Nor does approval solve representation errors. A mistaken equivalence,
misleading source mapping, fabricated receipt, or incomplete consequence
contract can lead a deterministic evaluator to an incorrect permit. The
system needs domain validation and monitoring in addition to formal
consistency. Shared meanings can propagate a mistake across many agents,
making governed semantic releases valuable and consequential at the same
time.

\section{Develop, operate, and
improve}\label{develop-operate-and-improve}

A shared semantic basis can support three functions with different
grants.

\textbf{Develop} implements capabilities and adapters against an
identified model and policy contract. Its permission to modify code or
author test data does not entail permission to release production
semantics, rewrite live mandates, or commit domain decisions.
Development output remains a proposed artifact until the release
authority admits it.

\textbf{Operate} interprets admitted situations and proposes or executes
decisions within the active contract. It can collect observations and
report mismatches. It cannot grant itself a broader decision class or
make a candidate ontology extension authoritative simply because the
extension appears useful.

\textbf{Improve} investigates failures and gaps, compares alternatives,
and proposes changes. It can have a broader research scope than the
operating function while retaining narrower authority over effects.
Discovery of a new category, a successful test, or a persuasive
diagnosis does not itself approve adoption.

An orchestrator or another authorized mechanism governs transitions
between these functions. They may be different agents, services, teams,
or roles in one platform; naming three agents does not establish
separation of duties if they share unrestricted credentials. The full
develop/operate/improve lifecycle is outside this paper's scope. Here it
shows why access to a common ontology need not imply a common authority
envelope.

\section{Discussion and future work}\label{discussion-and-future-work}

\subsection{Relationship to existing
mechanisms}\label{relationship-to-existing-mechanisms}

The architecture combines established semantic authorization, normative
modeling, ontology-change analysis, and version-bound enforcement
{[}1-3, 12-19, 22{]}. Its organizing question is: \emph{Does a candidate
semantic change alter the decisions eligible under a stated delegation
and comparison domain, and which role may ratify that effect?} The
comparison target makes the effect of a semantic change explicit, while
the admission procedure connects that effect to a designated authority.
The taxonomy example illustrates this relationship in a restricted
fragment; it does not establish a general preservation theorem.

A typed schema, curated decision table, state machine, or existing
policy data model may represent a narrow domain more cheaply than an
ontology. The ontology proposal is appropriate when shared relations and
inference make policy interpretation depend on evolving meaning. It
should be compared with simpler representations rather than treated as a
mandatory foundation for every agent.

\subsection{Technical and institutional
limits}\label{technical-and-institutional-limits}

The formal model is relative to represented proposals, evidence, and
consequences. It cannot certify that the institution's mandate is
legitimate, that classifications are fair, or that the admitted evidence
is true. Classical ontology consistency and SHACL conformance do not
establish these claims.

Complete mediation may be infeasible across legacy tools and external
workflows. Direct credentials, unmodeled notification effects, and data
disclosed during retrieval can bypass a final write gate. Multiagent
composition creates additional problems: individually allowed proposals
can combine into an unauthorized outcome. Trace and aggregate
constraints need state, concurrency control, and clearly attributed
effects.

The semantic comparison problem can be expensive or undecidable for
sufficiently expressive languages and open domains. Extracted
dependencies can be incomplete. A narrower sound approximation may
refuse benign changes; a broader heuristic may miss harmful ones. An
indeterminate result must remain visible, with a policy defining which
authority may accept the residual uncertainty.

Institutions can also approve an overbroad grant. Separating approval
roles supports accountability; it does not guarantee prudent judgment.
Evidence disputes, appeal, jurisdiction, and override procedures require
domain-specific governance that this paper does not supply.

\subsection{What an implementation study should
test}\label{what-an-implementation-study-should-test}

A future study should compare four configurations on the same workflows
and evidence: resource/operation permission checks alone; an explicit
business-decision policy gate; that gate with version-bound semantics
and current grants; and that gate with authority-relative
semantic-change review. Baselines must be configured competently, since
otherwise the study merely compares careful modeling with a deliberately
weak system.

Tests should include subclass additions, equivalence and mapping
changes, ordinary new instances, disputed type assertions, imported
rules, reasoner changes, unrelated ontology edits, expired grants,
concurrent releases, stale evidence, bypass credentials, cumulative
credits, and recommendations that downstream systems finalize. Report
expansion witnesses, unauthorized effects, false refusals, review
workload, latency, and recovery separately. A bounded executable example
does not constitute an experiment on agent reliability or a deployment
evaluation.

Future work should specify a decidable comparison fragment, implement a
change assessor, validate the mandates, and compare detection accuracy
and cost against strong existing semantic-policy and ontology-change
controls. These comparisons should establish whether the additional
review detects consequential authority expansions at acceptable cost.

\section{Conclusion}\label{conclusion}

Delegated intelligence needs an explicit account of which judgments and
consequences an institution has entrusted to it. An independently
governed shared ontology can help express that account, while ontology
revisions, mappings, admitted assertions, and inference choices can also
change its effective scope. A policy can remain textually identical as
the decisions it permits expand.

Bounded agency, as developed here, separates acquiring or proposing
knowledge from adopting an authoritative operational basis and renewing
affected decision rights. The proposed architecture pairs
authority-relative semantic comparison with explicit ratification,
current version-bound authorization, and mediated effects. Future
evaluation should test whether this obligation detects consequential
semantic changes beyond existing controls and measure review burden and
runtime cost.

\section*{Declarations}\label{declarations}
\addcontentsline{toc}{section}{Declarations}

\textbf{Nature of the work.} This is a conceptual design study. The
cases and finite taxonomy illustration are synthetic thought
experiments. No human-subject study, production deployment, or
agent-performance experiment was conducted.

\textbf{Research artifacts.} Manuscript sources, the executable
finite-taxonomy illustration, its recorded output, and an example
change-admission schema are available in the supplementary archive and
\href{https://github.com/prosessleder/bounded-agency}{project
repository}.

\textbf{AI assistance.} OpenAI Codex assisted literature discovery,
manuscript drafting, generation of illustrative code, and internal
critical review. The author is responsible for the manuscript's content,
references, and conclusions.

\begin{thebibliography}{99}
\bibitem{ref1}
Andrzej Uszok, Jeffrey M. Bradshaw, Matthew Johnson, Renia Jeffers,
Austin Tate, Jeff Dalton, and Stuart Aitken. 2004. ``KAoS Policy
Management for Semantic Web Services.'' \emph{IEEE Intelligent Systems}
19(4): 32--41.
\href{https://doi.org/10.1109/MIS.2004.31}{doi:10.1109/MIS.2004.31}.

\bibitem{ref2}
Lalana Kagal, Tim Finin, and Anupam Joshi. 2003. ``A Policy Based
Approach to Security for the Semantic Web.'' \emph{ISWC}, LNCS 2870:
402--418.
\href{https://doi.org/10.1007/978-3-540-39718-2_26}{doi:10.1007/978-3-540-39718-2\_26}.

\bibitem{ref3}
Guido Boella and Leendert van der Torre. 2004. ``Regulative and
Constitutive Norms in Normative Multiagent Systems.'' \emph{KR 2004}:
255--265. \href{https://cdn.aaai.org/KR/2004/KR04-028.pdf}{Primary AAAI
paper}.

\bibitem{ref4}
Vincent C. Hu, David Ferraiolo, D. Richard Kuhn, Adam Schnitzer, Kenneth
Sandlin, Robert Miller, and Karen Scarfone. 2014, updated 2019.
\emph{Guide to Attribute Based Access Control (ABAC) Definition and
Considerations}. NIST SP 800-162.
\href{https://doi.org/10.6028/NIST.SP.800-162}{doi:10.6028/NIST.SP.800-162}.

\bibitem{ref5}
Ruoming Pang, Ramon Caceres, Mike Burrows, Zhifeng Chen, Pratik Dave,
Nathan Germer, Alexander Golynski, Kevin Graney, Nina Kang, Lea Kissner,
Jeffrey L. Korn, Abhishek Parmar, Christina D. Richards, and Mengzhi
Wang. 2019. ``Zanzibar: Google's Consistent, Global Authorization
System.'' \emph{USENIX ATC}: 33--46.
\href{https://www.usenix.org/conference/atc19/presentation/pang}{Primary
USENIX paper and record}.

\bibitem{ref6}
Jaehong Park and Ravi Sandhu. 2004. ``The UCONABC Usage Control Model.''
\emph{ACM Transactions on Information and System Security} 7(1):
128--174.
\href{https://doi.org/10.1145/984334.984339}{doi:10.1145/984334.984339}.

\bibitem{ref7}
Joseph W. Cutler, Craig Disselkoen, Aaron Eline, Shaobo He, Kyle
Headley, Michael Hicks, Kesha Hietala, Eleftherios Ioannidis, John
Kastner, Anwar Mamat, Darin McAdams, Matt McCutchen, Neha Rungta, Emina
Torlak, and Andrew Wells. 2024. ``Cedar: A New Language for Expressive,
Fast, Safe, and Analyzable Authorization (Extended Version).''
arXiv:2403.04651v2. \href{https://arxiv.org/abs/2403.04651v2}{Primary
record}.

\bibitem{ref8}
Open Policy Agent project. Undated. \emph{Open Policy Agent
Documentation}. \href{https://www.openpolicyagent.org/docs}{Official
documentation}. Accessed 9 October 2026.

\bibitem{ref9}
Renato Iannella and Serena Villata, eds.~2018. \emph{ODRL Information
Model 2.2}. W3C Recommendation, 15 February.
\href{https://www.w3.org/TR/2018/REC-odrl-model-20180215/}{Dated
specification}.

\bibitem{ref10}
Pascal Hitzler, Markus Krötzsch, Bijan Parsia, Peter F. Patel-Schneider,
and Sebastian Rudolph, eds.~2012. \emph{OWL 2 Web Ontology Language
Primer (Second Edition)}. W3C Recommendation, 11 December.
\href{https://www.w3.org/TR/2012/REC-owl2-primer-20121211/}{Dated
specification}.

\bibitem{ref11}
Holger Knublauch and Dimitris Kontokostas, eds.~2017. \emph{Shapes
Constraint Language (SHACL)}. W3C Recommendation, 20 July.
\href{https://www.w3.org/TR/2017/REC-shacl-20170720/}{Dated
specification}.

\bibitem{ref12}
Elena Botoeva, Boris Konev, Carsten Lutz, Vladislav Ryzhikov, Frank
Wolter, and Michael Zakharyaschev. 2017. ``Inseparability and
Conservative Extensions of Description Logic Ontologies: A Survey.''
\emph{Reasoning Web: Logical Foundation of Knowledge Graph Construction
and Query Answering}, LNCS 9885: 27--89. RW2016 event; publisher
publication 2017.
\href{https://doi.org/10.1007/978-3-319-49493-7_2}{doi:10.1007/978-3-319-49493-7\_2};
\href{https://arxiv.org/abs/1804.07805}{author-deposited text}.

\bibitem{ref13}
Sven Peldszus, Jens Bürger, Timo Kehrer, and Jan Jürjens. 2021.
``Ontology-driven Evolution of Software Security.'' \emph{Data \&
Knowledge Engineering} 134: 101907.
\href{https://doi.org/10.1016/j.datak.2021.101907}{doi:10.1016/j.datak.2021.101907}.

\bibitem{ref14}
Bo Nørregaard Jørgensen and Zheng Grace Ma. 2026. ``Governing Agentic AI
in Enterprise Workflows: A Bounded-Autonomy Framework for Delegated
Authority and Controlled Execution.'' \emph{Information} 17(9): 923.
Published 20 September.
\href{https://doi.org/10.3390/info17090923}{doi:10.3390/info17090923}.

\bibitem{ref15}
Anuj Kaul, Qianlong Lan, and Pranay Gupta. 2026. ``Behavioral Governance
for Autonomous AI Agents: The AgentBound Framework.''
arXiv:2606.30970v2, revised 1 July.
\href{https://arxiv.org/abs/2606.30970v2}{Primary record}.

\bibitem{ref16}
Parviz Shariff. 2026. ``The Authority Resolution Framework: A
Five-Domain Ontology for Governing Who and What Decides, at Scale.''
arXiv:2608.15832v1. \href{https://arxiv.org/abs/2608.15832v1}{Primary
record}.

\bibitem{ref17}
Yingqi Zhang. 2026. ``Agent libOS: A Runtime Substrate for
Capability-Controlled Self-Evolving LLM Agents.'' arXiv:2606.03895v3,
revised 18 August. \href{https://arxiv.org/abs/2606.03895v3}{Primary
record}.

\bibitem{ref18}
Basavesh Ammanaghatta Shivakumar, Swarn Priya, and Peng Gao. 2026.
``AgentFlow: A Flow-Centric Policy Language and Framework for Securing
LLM Agent Systems.'' arXiv:2608.22868v1.
\href{https://arxiv.org/abs/2608.22868v1}{Primary record}.

\bibitem{ref19}
Jesus Salas. 2026. ``Authority at Commit Time: Reject-and-Rerun
Semantics for Governed Agentic Systems.'' arXiv:2609.31490v1.
\href{https://arxiv.org/abs/2609.31490v1}{Primary record}.

\bibitem{ref20}
Fred B. Schneider. 2000. ``Enforceable Security Policies.'' \emph{ACM
Transactions on Information and System Security} 3(1): 30--50.
\href{https://www.cs.cornell.edu/fbs/publications/EnfSecPols.pdf}{Author-hosted
primary paper}.

\bibitem{ref21}
Eugene F. Fama and Michael C. Jensen. 1983. ``Separation of Ownership
and Control.'' \emph{Journal of Law and Economics} 26(2): 301--325.
\href{https://doi.org/10.1086/467037}{doi:10.1086/467037}.

\bibitem{ref22}
Inferal. Undated. \emph{Inferal Governed Artifact Evolution Ontology}.
Unofficial 1.0 Editor's Draft.
\href{https://ontology.inferal.com/modules/artifact-evolution/spec/}{Maintained
specification}. Accessed 9 October 2026; no dated immutable publication
established.

\bibitem{ref23}
Andrzej Uszok, Jeffrey M. Bradshaw, James Lott, Maggie Breedy, Larry
Bunch, Paul Feltovich, Matthew Johnson, and Hyuckchul Jung. 2008. ``New
Developments in Ontology-Based Policy Management: Increasing the
Practicality and Comprehensiveness of KAoS.'' \emph{IEEE POLICY}:
145--152.
\href{https://doi.org/10.1109/POLICY.2008.47}{doi:10.1109/POLICY.2008.47};
\href{https://ontology.ihmc.us/Papers/uszok_kaos_policy08.pdf}{primary
manuscript}.

\bibitem{ref24}
Nuno Lopes, Sabrina Kirrane, Antoine Zimmermann, Axel Polleres, and
Alessandra Mileo. 2012. ``A Logic Programming Approach for Access
Control over RDF.'' \emph{ICLP Technical Communications}, LIPIcs 17:
381--392.
\href{https://doi.org/10.4230/LIPIcs.ICLP.2012.381}{doi:10.4230/LIPIcs.ICLP.2012.381}.

\bibitem{ref25}
Andrew J. I. Jones and Marek Sergot. 1996. ``A Formal Characterisation
of Institutionalised Power.'' \emph{Logic Journal of the IGPL} 4(3):
427--443.
\href{https://doi.org/10.1093/jigpal/4.3.427}{doi:10.1093/jigpal/4.3.427}.

\bibitem{ref26}
Tobin South, Samuele Marro, Thomas Hardjono, Robert Mahari, Cedric
Deslandes Whitney, Dazza Greenwood, Alan Chan, and Alex Pentland. 2025.
``Authenticated Delegation and Authorized AI Agents.''
arXiv:2501.09674v1. \href{https://arxiv.org/abs/2501.09674v1}{Primary
record}.

\bibitem{ref27}
Nenad Tomašev, Matija Franklin, and Simon Osindero. 2026. ``Intelligent
AI Delegation.'' arXiv:2602.11865v1.
\href{https://arxiv.org/abs/2602.11865v1}{Primary record}.

\bibitem{ref28}
Lalana Kagal. Undated. \emph{Rei Ontology Specifications, Ver 2.0}.
\href{https://userpages.cs.umbc.edu/lkagal1/rei/}{Author-maintained
specification}. Accessed 9 October 2026.
\end{thebibliography}
\end{document}
